% GENERATED FILE. Edit canonical lessons, then run npm run generate:books.
% canonical-source-hash: fnv1a64:a0e1952594d214aa
% canonical-lessons: PT-C323-consoada, PT-C323-presepio, PT-C323-passagem-de-ano, PT-C323-folar, PT-C323-amendoas

\chapter{A consoada}
\label{ch:pt-a-consoada}
\hlchaptermodality{323}

\begin{chapteropening}
\textbf{By the end of this chapter:} \emph{I can describe family feasts of the year: Christmas Eve dinner, the nativity scene, New Year’s Eve and Easter treats.}

Complete the last lesson of chapter 323: I can describe family feasts of the year: Christmas Eve dinner, the nativity scene, New Year’s Eve and Easter treats.
\end{chapteropening}

\section[a consoada]{a consoada — the Christmas Eve dinner}

\label{lesson:PT-C323-consoada}



\begin{quote}
Before the new one: say the Portuguese for to decorate, then the Portuguese for the decoration.
\end{quote}

\subsection*{You'll want to know: a consoada}
\textbf{a consoada} — \textquotedblleft{}the Christmas Eve dinner\textquotedblright{}.

\begin{quote}
Na consoada comemos bacalhau com batatas.
\end{quote}

\begin{quote}
\emph{At Christmas Eve dinner we eat salt cod with potatoes.}
\end{quote}

A Portuguese tradition: the family supper on the night of 24 December.

\begin{cousinweb}
\emph{consoada} is built on \emph{consoar} \textquotedblleft{}to have a light supper\textquotedblright{}, which is usually traced to Latin \emph{cōnsōlārī} \textquotedblleft{}to comfort\textquotedblright{} (Lewis \& Short); the path is not fully settled.
\end{cousinweb}

\begin{grammarlens}[title={What you've built}]
One more word for telling, explaining and describing at B1.
\end{grammarlens}

\subsection*{Your turn}
\begin{itemize}
\raggedright
  \item \emph{Say it:} \emph{a consoada}
  \item \emph{Say it:} \emph{a consoada}, once more
  \item \emph{Say it:} a sentence of your own with \emph{a consoada}
  \item \emph{Recall:} say the Portuguese for to decorate, then the Portuguese for the decoration, then say \emph{a consoada} again
\end{itemize}

\begin{tcolorbox}[breakable,colback=teal!4,colframe=teal!35!black,title={Before you move on}]
What does a consoada mean? (\textquotedblleft{}The Christmas Eve dinner\textquotedblright{}.) Say it once more.
\end{tcolorbox}
\section[o presépio]{o presépio — the nativity scene, the crib}

\label{lesson:PT-C323-presepio}



\begin{quote}
Before the new one: say the Portuguese for the decoration, then the Portuguese for the Christmas Eve dinner.
\end{quote}

\subsection*{You'll want to know: o presépio}
\textbf{o presépio} — \textquotedblleft{}the nativity scene, the crib\textquotedblright{}.

\begin{quote}
Todos os anos, a minha avó monta o presépio na sala.
\end{quote}

\begin{quote}
\emph{Every year my grandmother sets up the nativity scene in the living room.}
\end{quote}

\emph{Montar o presépio} — to set up the crib.

\begin{cousinweb}
\emph{presépio} is a learned borrowing from Latin \emph{praesaepium} \textquotedblleft{}a stall, a manger\textquotedblright{} (EtymDB).
\end{cousinweb}

\begin{grammarlens}[title={What you've built}]
One more word for telling, explaining and describing at B1.
\end{grammarlens}

\subsection*{Your turn}
\begin{itemize}
\raggedright
  \item \emph{Say it:} \emph{o presépio}
  \item \emph{Say it:} \emph{o presépio}, once more
  \item \emph{Say it:} a sentence of your own with \emph{o presépio}
  \item \emph{Recall:} say the Portuguese for the decoration, then the Portuguese for the Christmas Eve dinner, then say \emph{o presépio} again
\end{itemize}

\begin{tcolorbox}[breakable,colback=teal!4,colframe=teal!35!black,title={Before you move on}]
What does o presépio mean? (\textquotedblleft{}The nativity scene, the crib\textquotedblright{}.) Say it once more.
\end{tcolorbox}
\section[a passagem de ano]{a passagem de ano — the New Year celebration, New Year’s Eve}

\label{lesson:PT-C323-passagem-de-ano}



\begin{quote}
Before the new one: say the Portuguese for the Christmas Eve dinner, then the Portuguese for the nativity scene.
\end{quote}

\subsection*{You'll want to know: a passagem de ano}
\textbf{a passagem de ano} — \textquotedblleft{}the New Year celebration, New Year’s Eve\textquotedblright{}.

\begin{quote}
Na passagem de ano, vimos o fogo de artifício.
\end{quote}

\begin{quote}
\emph{On New Year’s Eve we watched the fireworks.}
\end{quote}

\emph{Feliz Ano Novo!} — Happy New Year! In Brazil: \emph{o réveillon}.

\begin{cousinweb}
\emph{passagem} came from French \emph{passage} (EtymDB), and \emph{ano} is inherited from Latin \emph{annus} \textquotedblleft{}a year\textquotedblright{}: the passing from one year to the next.
\end{cousinweb}

\begin{grammarlens}[title={What you've built}]
One more word for telling, explaining and describing at B1.
\end{grammarlens}

\subsection*{Your turn}
\begin{itemize}
\raggedright
  \item \emph{Say it:} \emph{a passagem de ano}
  \item \emph{Say it:} \emph{a passagem de ano}, once more
  \item \emph{Say it:} a sentence of your own with \emph{a passagem de ano}
  \item \emph{Recall:} say the Portuguese for the Christmas Eve dinner, then the Portuguese for the nativity scene, then say \emph{a passagem de ano} again
\end{itemize}

\begin{tcolorbox}[breakable,colback=teal!4,colframe=teal!35!black,title={Before you move on}]
What does a passagem de ano mean? (\textquotedblleft{}The New Year celebration, New Year’s Eve\textquotedblright{}.) Say it once more.
\end{tcolorbox}
\section[o folar]{o folar — the Easter bread}

\label{lesson:PT-C323-folar}



\begin{quote}
Before the new one: say the Portuguese for the nativity scene, then the Portuguese for the New Year celebration.
\end{quote}

\subsection*{You'll want to know: o folar}
\textbf{o folar} — \textquotedblleft{}the Easter bread\textquotedblright{}.

\begin{quote}
A madrinha oferece o folar ao afilhado.
\end{quote}

\begin{quote}
\emph{The godmother gives the Easter bread to her godson.}
\end{quote}

A sweet or savoury loaf, sometimes with whole eggs baked in.

\begin{cousinweb}
The origin of \emph{folar} is not settled; EtymDB has the word with no source.
\end{cousinweb}

\begin{grammarlens}[title={What you've built}]
One more word for telling, explaining and describing at B1.
\end{grammarlens}

\subsection*{Your turn}
\begin{itemize}
\raggedright
  \item \emph{Say it:} \emph{o folar}
  \item \emph{Say it:} \emph{o folar}, once more
  \item \emph{Say it:} a sentence of your own with \emph{o folar}
  \item \emph{Recall:} say the Portuguese for the nativity scene, then the Portuguese for the New Year celebration, then say \emph{o folar} again
\end{itemize}

\begin{tcolorbox}[breakable,colback=teal!4,colframe=teal!35!black,title={Before you move on}]
What does o folar mean? (\textquotedblleft{}The Easter bread\textquotedblright{}.) Say it once more.
\end{tcolorbox}
\section[as amêndoas]{as amêndoas — the sugared almonds; almonds}

\label{lesson:PT-C323-amendoas}



\begin{quote}
Before the new one: say the Portuguese for the New Year celebration, then the Portuguese for the Easter bread.
\end{quote}

\subsection*{You'll want to know: as amêndoas}
\textbf{as amêndoas} — \textquotedblleft{}the sugared almonds; almonds\textquotedblright{}.

\begin{quote}
As crianças receberam amêndoas da madrinha.
\end{quote}

\begin{quote}
\emph{The children got sugared almonds from their godmother.}
\end{quote}

At Easter, \emph{as amêndoas} are sugared or chocolate almonds.

\begin{cousinweb}
\emph{amêndoa} goes back to Late Latin \emph{amyndala}, from Latin \emph{amygdala}, a borrowing from Greek \emph{amygdálē} \textquotedblleft{}an almond\textquotedblright{} (Lewis \& Short).
\end{cousinweb}

\begin{grammarlens}[title={What you've built}]
One more word for telling, explaining and describing at B1.
\end{grammarlens}

\subsection*{Your turn}
\begin{itemize}
\raggedright
  \item \emph{Say it:} \emph{as amêndoas}
  \item \emph{Say it:} \emph{as amêndoas}, once more
  \item \emph{Say it:} a sentence of your own with \emph{as amêndoas}
  \item \emph{Recall:} say the Portuguese for the New Year celebration, then the Portuguese for the Easter bread, then say \emph{as amêndoas} again
\end{itemize}

\begin{tcolorbox}[breakable,colback=teal!4,colframe=teal!35!black,title={Before you move on}]
What does as amêndoas mean? (\textquotedblleft{}The sugared almonds; almonds\textquotedblright{}.) Say it once more.
\end{tcolorbox}
